% GENERATED FILE. Edit canonical lessons, then run npm run generate:books.
% canonical-source-hash: fnv1a64:3af119418da98bad
% canonical-lessons: GU-C171-svad, GU-C171-cakhvu, GU-C171-kadvu, GU-C171-kharu, GU-C171-galyu

\chapter{Bitter, salty and sweet}
\label{ch:gu-bitter-salty-and-sweet}
\hlchaptermodality{171}

\begin{chapteropening}
\textbf{By the end of this chapter:} \emph{I can say how food tastes.}

Complete the last lesson of chapter 171: I can say how food tastes.
\end{chapteropening}

\section[\texorpdfstring{\gu{સ્વાદ} (svād)}{સ્વાદ (svād)}]{\gu{સ્વાદ} (svād) — taste, flavour}

\label{lesson:GU-C171-svad}



\begin{quote}
Before the new one: say the Gujarati for to sift, then the Gujarati for to stick.
\end{quote}

\subsection*{You'll want to know: \gu{સ્વાદ}}
\textbf{\gu{સ્વાદ}} — \emph{svād} — \textquotedblleft{}taste, flavour\textquotedblright{}.

\begin{quote}
\textbf{\gu{આ} \gu{દાળનો} \gu{સ્વાદ} \gu{સારો} \gu{છે}.} — \emph{ā dāḷno svād sāro che} — \textquotedblleft{}This dal tastes good.\textquotedblright{}
\end{quote}

\begin{cousinweb}
\textbf{\gu{સ્વાદ}} (\emph{svād}) is Sanskrit \emph{svāda}, \textquotedblleft{}taste, flavour\textquotedblright{}, taken over whole as a learned loan, a tatsama. CDIAL lists the Sanskrit word but no Gujarati form (CDIAL 13923).
\end{cousinweb}

\begin{grammarlens}[title={What you've built}]
One more word for everyday life at A2.
\end{grammarlens}

\subsection*{Your turn}
\begin{itemize}
\raggedright
  \item \emph{Say it:} \emph{svād}
  \item \emph{Say it:} \emph{svād}, once more
  \item \emph{Say it:} \emph{ā dāḷno svād sāro che} — This dal tastes good.
  \item \emph{Recall:} say the Gujarati for to sift, then the Gujarati for to stick, then say \emph{svād} again
\end{itemize}

\begin{tcolorbox}[breakable,colback=teal!4,colframe=teal!35!black,title={Before you move on}]
What does \gu{સ્વાદ} mean? (\textquotedblleft{}Taste, flavour\textquotedblright{}.) Say it once more.
\end{tcolorbox}
\section[\texorpdfstring{\gu{ચાખવું} (cākhvũ)}{ચાખવું (cākhvũ)}]{\gu{ચાખવું} (cākhvũ) — to taste}

\label{lesson:GU-C171-cakhvu}



\begin{quote}
Before the new one: say the Gujarati for to stick, then the Gujarati for taste.
\end{quote}

\subsection*{You'll want to know: \gu{ચાખવું}}
\textbf{\gu{ચાખવું}} — \emph{cākhvũ} — \textquotedblleft{}to taste\textquotedblright{}.

\begin{quote}
\textbf{\gu{માતા} \gu{શાક} \gu{ચાખે} \gu{છે}.} — \emph{mātā śāk cākhe che} — \textquotedblleft{}Mother tastes the curry.\textquotedblright{}
\end{quote}

\begin{cousinweb}
\textbf{\gu{ચાખવું}} (\emph{cākhvũ}) is a tadbhava. It comes from \emph{*cakṣati}, \textquotedblleft{}tastes\textquotedblright{}, a reconstructed form, and CDIAL glosses the Gujarati verb \textquotedblleft{}to taste\textquotedblright{} (CDIAL 4557).
\end{cousinweb}

\begin{grammarlens}[title={What you've built}]
One more word for everyday life at A2.
\end{grammarlens}

\subsection*{Your turn}
\begin{itemize}
\raggedright
  \item \emph{Say it:} \emph{cākhvũ}
  \item \emph{Say it:} \emph{cākhvũ}, once more
  \item \emph{Say it:} \emph{mātā śāk cākhe che} — Mother tastes the curry.
  \item \emph{Recall:} say the Gujarati for to stick, then the Gujarati for taste, then say \emph{cākhvũ} again
\end{itemize}

\begin{tcolorbox}[breakable,colback=teal!4,colframe=teal!35!black,title={Before you move on}]
What does \gu{ચાખવું} mean? (\textquotedblleft{}To taste\textquotedblright{}.) Say it once more.
\end{tcolorbox}
\section[\texorpdfstring{\gu{કડવું} (kaḍvũ)}{કડવું (kaḍvũ)}]{\gu{કડવું} (kaḍvũ) — bitter}

\label{lesson:GU-C171-kadvu}



\begin{quote}
Before the new one: say the Gujarati for taste, then the Gujarati for to taste.
\end{quote}

\subsection*{You'll want to know: \gu{કડવું}}
\textbf{\gu{કડવું}} — \emph{kaḍvũ} — \textquotedblleft{}bitter\textquotedblright{}.

\begin{quote}
\textbf{\gu{આ} \gu{દવા} \gu{કડવી} \gu{છે}.} — \emph{ā davā kaḍvī che} — \textquotedblleft{}This medicine is bitter.\textquotedblright{}
\end{quote}

\begin{cousinweb}
\textbf{\gu{કડવું}} (\emph{kaḍvũ}) is a tadbhava. It comes from Sanskrit \emph{kaṭu}, \emph{kaṭuka}, \textquotedblleft{}pungent, bitter\textquotedblright{}, and CDIAL lists the Gujarati \emph{kaḍvũ} (CDIAL 2641).
\end{cousinweb}

\begin{grammarlens}[title={What you've built}]
One more word for everyday life at A2.
\end{grammarlens}

\subsection*{Your turn}
\begin{itemize}
\raggedright
  \item \emph{Say it:} \emph{kaḍvũ}
  \item \emph{Say it:} \emph{kaḍvũ}, once more
  \item \emph{Say it:} \emph{ā davā kaḍvī che} — This medicine is bitter.
  \item \emph{Recall:} say the Gujarati for taste, then the Gujarati for to taste, then say \emph{kaḍvũ} again
\end{itemize}

\begin{tcolorbox}[breakable,colback=teal!4,colframe=teal!35!black,title={Before you move on}]
What does \gu{કડવું} mean? (\textquotedblleft{}Bitter\textquotedblright{}.) Say it once more.
\end{tcolorbox}
\section[\texorpdfstring{\gu{ખારું} (khārũ)}{ખારું (khārũ)}]{\gu{ખારું} (khārũ) — salty}

\label{lesson:GU-C171-kharu}



\begin{quote}
Before the new one: say the Gujarati for to taste, then the Gujarati for bitter.
\end{quote}

\subsection*{You'll want to know: \gu{ખારું}}
\textbf{\gu{ખારું}} — \emph{khārũ} — \textquotedblleft{}salty\textquotedblright{}.

\begin{quote}
\textbf{\gu{આજે} \gu{દાળ} \gu{ખારી} \gu{છે}.} — \emph{āje dāḷ khārī che} — \textquotedblleft{}Today the dal is salty.\textquotedblright{}
\end{quote}

\begin{cousinweb}
\textbf{\gu{ખારું}} (\emph{khārũ}) is a tadbhava. It comes from Sanskrit \emph{kṣāra}, \textquotedblleft{}corrosive; alkali\textquotedblright{}, and CDIAL glosses the Gujarati word \textquotedblleft{}saline\textquotedblright{} (CDIAL 3674).
\end{cousinweb}

\begin{grammarlens}[title={What you've built}]
One more word for everyday life at A2.
\end{grammarlens}

\subsection*{Your turn}
\begin{itemize}
\raggedright
  \item \emph{Say it:} \emph{khārũ}
  \item \emph{Say it:} \emph{khārũ}, once more
  \item \emph{Say it:} \emph{āje dāḷ khārī che} — Today the dal is salty.
  \item \emph{Recall:} say the Gujarati for to taste, then the Gujarati for bitter, then say \emph{khārũ} again
\end{itemize}

\begin{tcolorbox}[breakable,colback=teal!4,colframe=teal!35!black,title={Before you move on}]
What does \gu{ખારું} mean? (\textquotedblleft{}Salty\textquotedblright{}.) Say it once more.
\end{tcolorbox}
\section[\texorpdfstring{\gu{ગળ્યું} (gaḷyũ)}{ગળ્યું (gaḷyũ)}]{\gu{ગળ્યું} (gaḷyũ) — sweet}

\label{lesson:GU-C171-galyu}



\begin{quote}
Before the new one: say the Gujarati for bitter, then the Gujarati for salty.
\end{quote}

\subsection*{You'll want to know: \gu{ગળ્યું}}
\textbf{\gu{ગળ્યું}} — \emph{gaḷyũ} — \textquotedblleft{}sweet\textquotedblright{}.

\begin{quote}
\textbf{\gu{ગુજરાતી} \gu{દાળ} \gu{ગળી} \gu{છે}.} — \emph{gujarātī dāḷ gaḷī che} — \textquotedblleft{}Gujarati dal is sweet.\textquotedblright{}
\end{quote}

\begin{cousinweb}
\textbf{\gu{ગળ્યું}} (\emph{gaḷyũ}) is a tadbhava. CDIAL lists it, glossed \textquotedblleft{}sweet\textquotedblright{}, among the words built on Sanskrit \emph{guḍa}, \textquotedblleft{}boiled sugarcane juice, molasses\textquotedblright{} (CDIAL 4182).
\end{cousinweb}

\begin{grammarlens}[title={What you've built}]
One more word for everyday life at A2.
\end{grammarlens}

\subsection*{Your turn}
\begin{itemize}
\raggedright
  \item \emph{Say it:} \emph{gaḷyũ}
  \item \emph{Say it:} \emph{gaḷyũ}, once more
  \item \emph{Say it:} \emph{gujarātī dāḷ gaḷī che} — Gujarati dal is sweet.
  \item \emph{Recall:} say the Gujarati for bitter, then the Gujarati for salty, then say \emph{gaḷyũ} again
\end{itemize}

\begin{tcolorbox}[breakable,colback=teal!4,colframe=teal!35!black,title={Before you move on}]
What does \gu{ગળ્યું} mean? (\textquotedblleft{}Sweet\textquotedblright{}.) Say it once more.
\end{tcolorbox}
